%2multibyte Version: 5.50.0.2960 CodePage: 1252

\documentclass{article}
%%%%%%%%%%%%%%%%%%%%%%%%%%%%%%%%%%%%%%%%%%%%%%%%%%%%%%%%%%%%%%%%%%%%%%%%%%%%%%%%%%%%%%%%%%%%%%%%%%%%%%%%%%%%%%%%%%%%%%%%%%%%%%%%%%%%%%%%%%%%%%%%%%%%%%%%%%%%%%%%%%%%%%%%%%%%%%%%%%%%%%%%%%%%%%%%%%%%%%%%%%%%%%%%%%%%%%%%%%%%%%%%%%%%%%%%%%%%%%%%%%%%%%%%%%%%
\usepackage{amsmath}

\setcounter{MaxMatrixCols}{10}
%TCIDATA{OutputFilter=LATEX.DLL}
%TCIDATA{Version=5.50.0.2960}
%TCIDATA{Codepage=1252}
%TCIDATA{<META NAME="SaveForMode" CONTENT="1">}
%TCIDATA{BibliographyScheme=Manual}
%TCIDATA{Created=Monday, September 22, 2003 09:40:11}
%TCIDATA{LastRevised=Wednesday, September 03, 2025 08:48:02}
%TCIDATA{<META NAME="GraphicsSave" CONTENT="32">}
%TCIDATA{<META NAME="DocumentShell" CONTENT="General\SW\Blank - Standard LaTeX Article">}
%TCIDATA{Language=American English}
%TCIDATA{CSTFile=LaTeX article (bright).cst}

\newtheorem{theorem}{Theorem}
\newtheorem{acknowledgement}[theorem]{Acknowledgement}
\newtheorem{algorithm}[theorem]{Algorithm}
\newtheorem{axiom}[theorem]{Axiom}
\newtheorem{case}[theorem]{Case}
\newtheorem{claim}[theorem]{Claim}
\newtheorem{conclusion}[theorem]{Conclusion}
\newtheorem{condition}[theorem]{Condition}
\newtheorem{conjecture}[theorem]{Conjecture}
\newtheorem{corollary}[theorem]{Corollary}
\newtheorem{criterion}[theorem]{Criterion}
\newtheorem{definition}[theorem]{Definition}
\newtheorem{example}[theorem]{Example}
\newtheorem{exercise}[theorem]{Exercise}
\newtheorem{lemma}[theorem]{Lemma}
\newtheorem{notation}[theorem]{Notation}
\newtheorem{problem}[theorem]{Problem}
\newtheorem{proposition}[theorem]{Proposition}
\newtheorem{remark}[theorem]{Remark}
\newtheorem{solution}[theorem]{Solution}
\newtheorem{summary}[theorem]{Summary}
\newenvironment{proof}[1][Proof]{\textbf{#1.} }{\ \rule{0.5em}{0.5em}}
\input{tcilatex}
\input{tcilatex}
\setlength{\textwidth}{6.5in}
\setlength{\textheight}{9in}
\setlength{\topmargin}{-.50in}
\setlength{\oddsidemargin}{.0in}
\setlength{\evensidemargin}{.0in}

\begin{document}


\begin{center}
{\LARGE Assignment 2}
\end{center}

\noindent Blankenau\hspace*{\stretch{0.7500}}

\noindent Economics 905, Fall 2025\medskip

\noindent \textit{Instructions.} You may work in groups and may compare
answers prior to submission. Joint submissions are allowed and encouraged.
All work should be legible and concise. A due date will be announced in
class. \ \ \ \ \ \ \ \ \ \ \ \ \ \ \ \ \ \ \ \ \ \ \ \ \ \ \ \ \ \ \ \ \ \ \
\ \ \ \ \ \ \ \ \ \ \ \ \ \ \ \ \ \ \ \ \ \ \ \ \ \ \ \ \ \ \ 

\textbf{Note: }I have included the key here. You are to try this as best you
can and look at the key when you can no longer make progress. Despite having
the key, you are expected to turn in the assignment.

\begin{enumerate}
\item Consider a static economy with a single representative agent endowed
with one unit of time to allocate across labor and leisure and $k_{0}$ units
of capital. Preference are given by 
\begin{equation}
U=\ln c+\beta \ln \ell  \label{utility}
\end{equation}%
where $\beta >0,$ $c$ is the agent's consumption and $\ell $ is the agent's
leisure$.$ There is also a representative firm with a production technology
given by 
\begin{equation*}
y=k^{\alpha }n^{1-\alpha }.
\end{equation*}%
where $\alpha \in \left( 0,1\right) ,$ $k$ is capital employed by the firm
and $n$ is the labor employed by the firm$.$ Finally there is a government
which consumes $G$ units of output, finances its consumption by taxing and
is subject to a balanced budget constraint.

\begin{enumerate}
\item State the social planner's problem and find the first order conditions
for the social planner's problem.

\item Find an implicit function which could be used to solve for $n.$ This
will be non-linear and you will not be able to solve explicitly for $n.$

\item Now lets look at the competitive equilibrium assuming that the
government pays for its expenditures through a lump sum tax. First, state
the agent's problem and show that optimization requires 
\begin{equation*}
\frac{w}{wn_{s}+rk_{s}-\tau }=\frac{\beta }{1-n_{s}}
\end{equation*}%
where $r$ is the rental rate of capital, $w$ is the wage rate for a unit of
labor and $\tau $ is the lump sum tax. The $s$ subscript denotes quantities
supplied.

\item State the firm's problem. From the first order conditions, find
expressions for $r$ and $w$.

\item Find the government's balanced budget constraint.

\item Define a competitive equilibrium in this case.

\item Use the information from d-f to arrive at an equation identical to
that in b. Why does this occur?

\item Now suppose an identical environment except that taxes are imposed at
an equal rate on all income. That is, the tax burden on an individual is $%
\hat{\tau}\left( wn_{s}+rk_{s}\right) $ and net income is $\left( 1-\hat{\tau%
}\right) \left( wn_{s}+rk_{s}\right) .$ Show that this time the solution to
the social planner's problem and the solution that arises in a competitive
equilibrium differ. You will not need to solve the social planner's problem
again. (It would be good to think about why but that is on your own). Thus
you need to show that the competitive equilibrium gives something else. Why
is this?

\item Now lets consider an identical environment except that the government
collects a fixed share of output rather than a fixed quantity of output. 
\textit{It again finances consumption through a lump sum tax. }That is, set $%
G=gy$ so that the size of government depends on the size of output. Solve
the social planner's problem for the optimal amount of $n$. The social
planner does not choose $g$. This time you will be able to solve explicitly
for the optimal $n$.

\item Solve the competitive equilibrium in this case. Note that the agent
takes his lump sum tax as given. That is, since he is a representative
agent, he thinks of himself as one agent of many. He knows that if total
income goes up, so will government spending. However, an increase in his own
income is too small to have a measurable effect on total income. Thus the
agent's problem is the same as before. It is the government's budget
constraint that is different. You will be able to solve explicitly for the
equilibrium level of $n$.

\item Compare your answers to the previous two parts and comment. Are they
the same or different? Why?

\item Show that the utility of an agent under the social planner's solution
and in a competitive equilibrium can be written respectively as 
\begin{equation}
U=\ln \left( 1-g\right) k^{\alpha }+\left( 1-\alpha \right) \ln \left( \frac{%
1-\alpha }{1+\beta -\alpha }\right) +\beta \ln \left( \frac{\beta }{1+\beta
-\alpha }\right)  \label{a}
\end{equation}%
and%
\begin{equation}
U=\ln \left( 1-g\right) k^{\alpha }+\left( 1-\alpha \right) \ln \left( \frac{%
1-\alpha }{1+\beta \left( 1-g\right) -\alpha }\right) +\beta \ln \left( 
\frac{\beta \left( 1-g\right) }{1+\beta \left( 1-g\right) -\alpha }\right)
\label{b}
\end{equation}

\item We want to measure the welfare cost of this inefficiency. Toward this,
consider scaling an agent's consumption of the good and leisure in the
competitive equilibrium so that in a competitive equilibrium 
\begin{equation*}
U=\ln \left( zc\right) +\beta \ln \left( z\ell \right) .
\end{equation*}%
Show that if 
\begin{equation*}
z=\frac{1}{\left( 1-g\right) ^{\frac{\beta }{1+\beta }}}\left( \frac{1+\beta
\left( 1-g\right) -\alpha }{1+\beta -\alpha }\right) ^{\frac{1-\alpha +\beta 
}{1+\beta }}
\end{equation*}%
the utility in the competitive equilibrium is the same as in the social
planner's problem. That is, you want to equate \ref{a} and your new version
of \ref{b}, solve for $z$, and show that it reduces to the above expression$%
. $

\item Interpret $z.$
\end{enumerate}

\pagebreak
\end{enumerate}

\begin{enumerate}
\item ..

\begin{enumerate}
\item Note that in this solution I am assuming interior solutions. In
general we should be careful about this but with log preferences we know
that neither $c$ nor \ $\ell $ will be 0. I skip a bit of this and jump into
the Lagrangian for the social planner's problem:%
\begin{equation*}
L=\ln c+\beta \ln \ell +\lambda _{1}\left( k^{\alpha }n^{1-\alpha
}-c-G\right) +\lambda _{2}\left( 1-\ell -n\right)
\end{equation*}%
which simplifies to%
\begin{equation*}
L=\ln (k_{0}^{\alpha }n^{1-\alpha }-G)+\beta \ln \left( 1-n\right)
\end{equation*}

\item The FOC gives (after some simplification):%
\begin{equation}
\frac{\left( 1-\alpha \right) k_{0}^{\alpha }n^{-\alpha }}{k_{0}^{\alpha
}n^{1-\alpha }-G}=\frac{\beta }{1-n}  \label{OPT1}
\end{equation}%
Now turn to the competitive equilibrium. \ 

\item Agent's problem%
\begin{equation*}
L=\ln c+\beta \ln \ell +\lambda _{1}\left( wn+rk_{0}-c-\tau \right)
\end{equation*}%
\begin{equation*}
L=\ln \left( wn+rk_{0}-\tau \right) +\beta \ln \left( 1-n\right)
\end{equation*}%
First order conditions give%
\begin{equation}
\frac{w}{wn+rk_{0}-\tau }=\frac{\beta }{1-n}  \label{foc}
\end{equation}

\item The firm's maximization gives 
\begin{equation}
r=\alpha \left( \frac{n}{k}\right) ^{1-\alpha }  \label{r}
\end{equation}%
\begin{equation}
w=\left( 1-\alpha \right) \left( \frac{k}{n}\right) ^{\alpha }  \label{w}
\end{equation}

\item Government budget constraint gives $G=\tau .$

\item A competitive equilibrium in this economy is a set of quantities $%
\left( c,\ell ,n,\tau ,G\right) $ and prices $\left( \omega ,r\right) $
which satisfy the following:

\begin{enumerate}
\item Consumer optimizes: equation (\ref{foc}).

\item Firm maximizes profits: equations (\ref{r}) and (\ref{w}).

\item The goods market, labor, market and capital markets clear; i.e. $%
y=c+G, $ capital demanded equals capital supplied and labor demanded equals
labor supplied.

\item The government budget constraint is satisfied; $G=\tau $.
\end{enumerate}

\item Put these and $k=k_{0}$ into \ref{foc} and simplify to get \ 
\begin{equation*}
\frac{\left( 1-\alpha \right) \left( \frac{k}{n}\right) ^{\alpha }}{%
k^{\alpha }n^{1-\alpha }-G}=\frac{\beta }{1-n}
\end{equation*}%
Inspect to see these are the same. While we cannot solve for $n$ explicitly
it is obvious that the same $n$ will solve both the competitive equilibrium
and the social planner's problem. This of course holds because we are
satisfying the requirements of the welfare theorems.

\item If the tax is instead on income the agent's problem becomes 
\begin{equation*}
L=\ln \left( \left( 1-\hat{\tau}\right) \left( wn+rk_{0}\right) \right)
+\beta \ln \left( 1-n\right)
\end{equation*}%
and the first order condition is 
\begin{equation*}
\frac{w}{wn+rk_{0}}=\frac{\beta }{1-n}
\end{equation*}%
There is nothing different about the firm's problem so put in wages and
rental rate of capital to get 
\begin{equation*}
\frac{\left( 1-\alpha \right) }{n}=\frac{\beta }{1-n}
\end{equation*}%
Clearly a different $n$ will solve this. This is not surprising since with a
distortionary tax, welfare theorems are not met and a competitive
equilibrium is not the same as the solution to the social planner's problem.

\item When the government collects a fixed share of output the social
planner's problem becomes%
\begin{equation*}
L=\ln c+\beta \ln \ell +\lambda _{1}\left( \left( 1-g\right) k^{\alpha
}n^{1-\alpha }-c\right) +\lambda _{2}\left( 1-\ell -n\right)
\end{equation*}%
or%
\begin{equation*}
L=\ln (\left( 1-g\right) k^{\alpha }n^{1-\alpha })+\beta \ln \left(
1-n\right)
\end{equation*}%
First order conditions give 
\begin{equation*}
\frac{\left( 1-g\right) \left( 1-\alpha \right) k^{\alpha }n^{-\alpha }}{%
\left( 1-g\right) k^{\alpha }n^{1-\alpha }}=\frac{\beta }{1-n}
\end{equation*}%
or after solving for $n$ 
\begin{equation*}
n=\frac{1-\alpha }{1+\beta -\alpha }.
\end{equation*}

\item Now lets solve the competitive equilibrium. The agent's problem is the
same as before and so is the firm's so we know%
\begin{equation*}
\frac{\left( 1-\alpha \right) k_{0}^{\alpha }n^{-\alpha }}{k_{0}^{\alpha
}n^{1-\alpha }-G}=\frac{\beta }{1-n}
\end{equation*}%
but now $G=gy$ so we get 
\begin{equation*}
\frac{\left( 1-\alpha \right) k_{0}^{\alpha }n^{-\alpha }}{\left( 1-g\right)
k_{0}^{\alpha }n^{1-\alpha }}=\frac{\beta }{1-n}
\end{equation*}%
and simplifying gives 
\begin{equation*}
\frac{\left( 1-\alpha \right) }{\left( 1-g\right) n}=\frac{\beta }{1-n}
\end{equation*}%
or 
\begin{equation*}
\frac{\left( 1-\alpha \right) }{1+\beta \left( 1-g\right) -\alpha }=n
\end{equation*}

\item From this it is clear that in the competitive equilibrium agent's work
too much. Why? They ignore that when they work, it increases the tax burden.
For each agent this effect is negligible but taken together it leads to a
distortion. The social planner takes this into account.

\item Utility is%
\begin{equation*}
U=\ln c+\beta \ln \ell
\end{equation*}%
In the social planner's problem this is 
\begin{equation*}
U=\ln \left( \left( 1-g\right) k^{\alpha }\left( \frac{1-\alpha }{1+\beta
-\alpha }\right) ^{1-\alpha }\right) +\beta \ln \left( 1-\frac{1-\alpha }{%
1+\beta -\alpha }\right)
\end{equation*}%
\begin{equation*}
U=\ln \left( 1-g\right) k^{\alpha }+\left( 1-\alpha \right) \ln \left( \frac{%
1-\alpha }{1+\beta -\alpha }\right) +\beta \ln \left( \frac{\beta }{1+\beta
-\alpha }\right)
\end{equation*}%
whereas in a competitive equilibrium.%
\begin{equation*}
U=\ln \left( \left( 1-g\right) k^{\alpha }\left( \frac{\left( 1-\alpha
\right) }{1+\beta \left( 1-g\right) -\alpha }\right) ^{1-\alpha }\right)
+\beta \ln \left( 1-\frac{\left( 1-\alpha \right) }{1+\beta \left(
1-g\right) -\alpha }\right)
\end{equation*}%
\begin{equation*}
U=\ln \left( 1-g\right) k^{\alpha }+\left( 1-\alpha \right) \ln \left( \frac{%
1-\alpha }{1+\beta \left( 1-g\right) -\alpha }\right) +\beta \ln \left( 
\frac{\beta \left( 1-g\right) }{1+\beta \left( 1-g\right) -\alpha }\right)
\end{equation*}

\item Now scale all by $z$. This gives 
\begin{equation*}
U=\ln \left( z\left( 1-g\right) k^{\alpha }\left( \frac{\left( 1-\alpha
\right) }{1+\beta \left( 1-g\right) -\alpha }\right) ^{1-\alpha }\right)
+\beta \ln \left( z\left( 1-\frac{\left( 1-\alpha \right) }{1+\beta \left(
1-g\right) -\alpha }\right) \right)
\end{equation*}%
or 
\begin{equation*}
U=\ln \left( 1-g\right) k^{\alpha }+\left( 1-\alpha \right) \ln \left( \frac{%
1-\alpha }{1+\beta \left( 1-g\right) -\alpha }\right) +\beta \ln \left( 
\frac{\beta \left( 1-g\right) }{1+\beta \left( 1-g\right) -\alpha }\right)
+\left( 1+\beta \right) \ln z
\end{equation*}%
Now solve to set these equal 
\begin{eqnarray*}
&&\ln \left( 1-g\right) k^{\alpha }+\left( 1-\alpha \right) \ln \left( \frac{%
1-\alpha }{1+\beta \left( 1-g\right) -\alpha }\right) +\beta \ln \left( 
\frac{\beta \left( 1-g\right) }{1+\beta \left( 1-g\right) -\alpha }\right)
+\left( 1+\beta \right) \ln z \\
&=&\ln \left( 1-g\right) k^{\alpha }+\left( 1-\alpha \right) \ln \left( 
\frac{1-\alpha }{1+\beta -\alpha }\right) +\beta \ln \left( \frac{\beta }{%
1+\beta -\alpha }\right)
\end{eqnarray*}%
or%
\begin{eqnarray*}
&&\ln \left( \left( \frac{1}{1+\beta \left( 1-g\right) -\alpha }\right)
^{\left( 1-\alpha \right) }\left( \frac{\beta \left( 1-g\right) }{1+\beta
\left( 1-g\right) -\alpha }\right) ^{\beta }z^{\left( 1+\beta \right)
}\right) \\
&=&\ln \left( \left( \frac{1}{1+\beta -\alpha }\right) ^{\left( 1-\alpha
\right) }\left( \frac{\beta }{1+\beta -\alpha }\right) ^{\beta }\right)
\end{eqnarray*}%
or%
\begin{eqnarray*}
&&\left( \frac{1}{1+\beta \left( 1-g\right) -\alpha }\right) ^{\left(
1-\alpha \right) }\left( \frac{\beta \left( 1-g\right) }{1+\beta \left(
1-g\right) -\alpha }\right) ^{\beta }z^{\left( 1+\beta \right) } \\
&=&\left( \frac{1}{1+\beta -\alpha }\right) ^{\left( 1-\alpha \right)
}\left( \frac{\beta }{1+\beta -\alpha }\right) ^{\beta }
\end{eqnarray*}%
or%
\begin{equation*}
z=\frac{1}{\left( 1-g\right) ^{\frac{\beta }{1+\beta }}}\left( \frac{1+\beta
\left( 1-g\right) -\alpha }{1+\beta -\alpha }\right) ^{\frac{1-\alpha +\beta 
}{1+\beta }}
\end{equation*}

\item This is a measure of the welfare loss due to the inefficiency.
Consumption would have to be scaled by $z$ in each period to make the agent
as well off with the inefficiency as she would be without the inefficiency.
\end{enumerate}
\end{enumerate}

\end{document}
